Before presenting the findings, we analyze the real-world results from three aspects: the gap to human teleoperation, the mismatch between simulation and real-world rankings, and deployment-specific failures such as action jitter and unsafe behaviors. These analyses show that RoboDojo-RealEval complements simulation by exposing physical-world bottlenecks that are difficult to capture in simulation alone.

\textcolor{Blue_1}{\textbf{Finding 1: Real-world deployment remains unreliable, with partial progress rarely translating into task completion.}}
Table~\ref{tab:robodojo_real_world_benchmark} shows that RoboDojo-RealEval remains highly challenging for current generalist manipulation policies. The best-performing policy, $\pi_{0.5}$, achieves only a $12.8\%$ overall success rate and a $22.9$ score across 18 real-world tasks, while human teleoperation reaches a $100.0\%$ success rate and a $100.0$ score on all tasks and embodiments. This gap shows that the tasks are feasible under the standardized protocol, yet current policies remain far from reliable physical deployment.

The higher scores relative to success rates further indicate that current policies can often make partial progress but fail to complete the full task. For example, $\pi_{0.5}$, InternVLA-A1, and GalaxeaVLA obtain scores of $22.9$, $12.0$, and $9.0$, respectively, but their success rates remain only $12.8\%$, $7.2\%$, and $4.4\%$. This suggests that real-world failures are driven not only by task misunderstanding, but also by execution bottlenecks such as precise alignment, contact handling, sub-step transition, and recovery from intermediate errors. Reliable real-world manipulation therefore requires closed-loop execution that can convert partial progress into final task success under physical uncertainty.

\textcolor{Blue_1}{\textbf{Finding 2: Simulation performance and real-world deployability are only partially aligned.}}
Since RoboDojo-RealEval currently evaluates a subset of the simulation leaderboard, we focus on policies evaluated in both settings. Among these overlapping policies, the two leaderboards show partial but not complete alignment. $\pi_{0.5}$ remains the strongest policy in RoboDojo-RealEval and is also among the leading methods in simulation, suggesting that simulation performance captures important aspects of general manipulation capability. However, the relative ordering changes for several policies: InternVLA-A1 and GalaxeaVLA achieve stronger real-world positions than their simulation rankings would suggest, while some policies with competitive simulation performance do not preserve the same advantage after physical deployment.

This mismatch should be interpreted carefully. RoboDojo is not designed as a paired sim-to-real transfer benchmark with one-to-one matched simulation and real-world task distributions. Instead, the two settings stress complementary aspects of generalist manipulation policies. Simulation provides scalable capability diagnosis across generalization, precision, long-horizon execution, memory, and open-semantic grounding, while RoboDojo-RealEval measures whether policies remain executable and reliable under physical deployment constraints.

Real-world evaluation introduces factors that are difficult to fully capture in simulation, including perception noise, camera and robot calibration errors, actuation latency, controller-specific dynamics, contact instability, and accumulated execution drift. These factors can cause policies with strong simulation performance to fail if their actions are jittery, poorly calibrated, insufficiently contact-aware, or unable to recover from small deviations. The partial alignment therefore supports the sim-and-real design of RoboDojo: simulation enables efficient large-scale capability analysis, while real-world evaluation exposes deployment-specific bottlenecks that cannot be inferred from simulation alone.

\textcolor{Blue_1}{\textbf{Finding 3: Real-world evaluation exposes execution instability and safety-critical behaviors beyond aggregate success metrics.}}
Real-world evaluation reveals deployment-specific failure modes that are not fully captured by aggregate success rates or simulation scores. Across several policies, we observe action jitter, oscillatory motions, repeated ineffective actions, contact instability, and occasional safety-critical behaviors. These failures indicate that physical deployment depends not only on high-level task understanding, but also on the temporal stability and safety of low-level action execution.

For example, DM0 frequently produces unstable control signals during deployment, leading to erratic motions that require close monitoring or safety intervention. This distinction is important: in simulation, unstable actions may only reduce task scores, whereas on physical robots they can introduce hardware risk, unsafe contacts, or unintended interactions with the environment. Therefore, final task success alone is insufficient to characterize real-world deployment quality.

These observations suggest that future generalist manipulation policies should optimize deployment-oriented properties beyond task completion, including action smoothness, contact-aware feedback, bounded control behavior, and recovery from off-trajectory states. RoboDojo-RealEval therefore provides complementary diagnostic value by revealing whether a policy is not only capable in principle, but also stable and safe enough for physical execution.